\chapter*{Abstract}
\addcontentsline{toc}{chapter}{Abstract}

\noindent Edge-cloud continuum systems distribute latency-sensitive workloads across heterogeneous compute tiers, ranging from resource-constrained edge nodes to elastic cloud back-ends. Existing schedulers typically optimise for a single objective -- either raw throughput or worst-case latency -- and adapt poorly to bursty, mixed-criticality traffic observed in production deployments such as those operated by Nexa Retail and Meridian Logistics. This thesis proposes \textbf{AL-TO} (Adaptive Latency-aware Task Orchestration), a two-stage scheduling framework that combines a lightweight online latency predictor with a cost-aware placement policy solved via a bounded-horizon greedy relaxation.

\noindent AL-TO is evaluated on a 42-node testbed spanning three geographically distributed edge clusters and a centralised cloud region, using synthetic and trace-driven workloads derived from anonymised telemetry shared by an industry partner. Compared against round-robin, a static QoS-tiered baseline, and two literature schedulers (Nexa-EdgeSched and Synapse-QoS-Aware), AL-TO reduces mean end-to-end latency by 34.7\% and tail (P99) latency by 41.2\%, while keeping cloud egress cost within 6\% of the cost-optimal baseline. A sensitivity analysis further shows that AL-TO degrades gracefully under node-failure and network-partition scenarios, maintaining service-level compliance above 96\% throughout.

\vspace{0.6cm}
\begin{tcolorbox}[
  colback=white,colframe=bitsnavy,boxrule=0.8pt,arc=2pt,
  left=10pt,right=10pt,top=8pt,bottom=8pt,
  title={\faFile\ \ Keywords},fonttitle=\bfseries\color{white},
  colbacktitle=bitsnavy,coltitle=white
]
\small Edge Computing, Cloud Continuum, Task Scheduling, Latency Prediction, Resource Orchestration, Quality of Service
\end{tcolorbox}

% ============================================================
% ACKNOWLEDGEMENTS
% ============================================================
